\documentclass{article}
\usepackage{amsmath}
\usepackage{graphicx}
\usepackage{array}
\begin{document}
\begin{table}[h]
\caption{GA com Correlação de Pearson para a Simulação 1.}
\label{tab 2}
\begin{tabular}{c|c|c|c|c}
\hline
\textbf{Método} & \textbf{$\gamma$} & \textbf{\begin{tabular}[c]{@{}c@{}}\emph{SNPs} selecionados\end{tabular}} & \textbf{\# \emph{SNPs} (V)} & \textbf{\begin{tabular}[c]{@{}c@{}}Correlação\\média\end{tabular}} \\
\hline
SNPs causais & - & \begin{tabular}[c]{@{}c@{}}\textbf{1}, \textbf{2}, \textbf{3}, \textbf{4}, \textbf{5}\\\textbf{6}, \textbf{7}, \textbf{8}\end{tabular} & - & - \\
\hline
SMS Linear & - & \begin{tabular}[c]{@{}c@{}}\textbf{1}, \textbf{2}, \textbf{3}, \textbf{4}, \textbf{5}\\\textbf{6}, \textbf{7}, \textbf{8}, 10, 11\\12, 51, 60, 71, 82\\83, 88\end{tabular} & 17(8) & 0.51 \\
\hline
SMS Radial & 0,001 & \begin{tabular}[c]{@{}c@{}}\textbf{1}, \textbf{2}, \textbf{4}, \textbf{5}, \textbf{7}\\\textbf{8}, 10, 12, 51, 83\end{tabular} & 18(6) & 0.51 \\
\hline
SMS Radial & 0,01 & \begin{tabular}[c]{@{}c@{}}\textbf{1}, \textbf{2}, \textbf{4}, \textbf{5}, \textbf{7}\\\textbf{8}, 46, 71, 83\end{tabular} & 10(6) & 0.58 \\
\hline
SMS Radial & 0,1 & \begin{tabular}[c]{@{}c@{}}\textbf{1}, \textbf{2}, \textbf{4}, \textbf{5}, \textbf{7}\\\textbf{8}, 46, 71, 83\end{tabular} & 9(6) & 0.65 \\
\hline
SMS Radial & 1 & \begin{tabular}[c]{@{}c@{}}\textbf{1}, \textbf{2}, \textbf{5}, \textbf{6}, \textbf{7}\\\textbf{8}, 51\end{tabular} & 7(6) & 0.59 \\
\hline
União & - & \begin{tabular}[c]{@{}c@{}}\textbf{1}, \textbf{2}, \textbf{3}, \textbf{4}, \textbf{5}\\\textbf{6}, \textbf{7}, \textbf{8}, 10, 11\\12, 46, 51, 60, 67\\68, 71, 77, 82, 83\\88\end{tabular} & 21(8) & - \\
\hline
Interseção & - & \begin{tabular}[c]{@{}c@{}}\textbf{1}, \textbf{2}, \textbf{5}, \textbf{7}, \textbf{8}\end{tabular} & 5(5) & - \\
\hline
Valor-p bruto & - & \begin{tabular}[c]{@{}c@{}}\textbf{1}, \textbf{2}, \textbf{3}, \textbf{4}, \textbf{5}\\\textbf{6}, \textbf{7}, \textbf{8}, 10, 12\\60, 71, 88\end{tabular} & 13(8) & - \\
\hline
Valor-p corrigido & - & \begin{tabular}[c]{@{}c@{}}\textbf{1}, \textbf{4}, \textbf{7}, \textbf{8}\end{tabular} & 4(4) & - \\
\hline
\end{tabular}
\end{table}
\end{document}
